\documentclass[conference]{IEEEtran}
\IEEEoverridecommandlockouts

% --- Packages ---
\usepackage{cite}
\usepackage{amsmath,amssymb,amsfonts}
\usepackage{algorithmic}
\usepackage{graphicx}
\usepackage{textcomp}
\usepackage{xcolor}
\usepackage{booktabs}
\usepackage{tabularx}
\usepackage{hyperref}
\usepackage{multirow}
\usepackage{float}
\usepackage{algorithm}
\usepackage{enumitem}

\begin{document}

\title{AgriSense AI: An Integrated Architecture for Crop Recommendation, Market Price Forecasting, and Pathology Detection}

\author{\IEEEauthorblockN{1\textsuperscript{st} Karthik Reddy}
\IEEEauthorblockA{\textit{B.Tech Computer Science \& Artificial Intelligence} \\
\textit{Rishihood University}\\
Roll No. 230035}
\and
\IEEEauthorblockN{2\textsuperscript{nd} Akash G}
\IEEEauthorblockA{\textit{B.Tech Computer Science \& Artificial Intelligence} \\
\textit{Rishihood University}\\
Roll No. 230098}
}

\maketitle

\begin{abstract}
Agricultural volatility and biotic stress remain pressing socioeconomic challenges in India, affecting the livelihoods of over 150 million farming households. This project presents AgriSense AI, an integrated machine learning pipeline addressing three interconnected domains: agronomic crop recommendation, 14-day market price forecasting, and computer-vision-based plant disease detection. The Crop Recommendation subsystem synthesizes over 45 engineered soil and climatic features to yield near-perfect LightGBM classification precision. For market intelligence, we construct a 14-day price forecasting apparatus utilizing scraped Agmarknet pricing and NASA POWER atmospheric data. A multi-variate Temporal Fusion Transformer (TFT) is deployed, establishing sophisticated probabilistic quantile bounds that effectively process external shocks, outperforming both traditional LightGBM boosting and Salesforce's zero-shot Moirai architecture. Finally, localized diagnostic capabilities are introduced via a Convolutional Neural Network (EfficientNet-B0), trained on the Augmented Plant Village dataset to classify 38 unique crop-disease intersections with profound accuracy. Ultimately, these three discrete deep learning and gradient-boosted subsystems are fused into a monolithic Python HTTP server, offering farmers an end-to-end deterministic dashboard for decision making traversing seed selection to harvest pathology and market realization.
\end{abstract}

\begin{IEEEkeywords}
Temporal Fusion Transformer, EfficientNet-B0, LightGBM, Agricultural Forecasting, Plant Pathology, Moirai, Computer Vision.
\end{IEEEkeywords}

\section{Introduction}

\subsection{Background and Motivation}
India's agricultural sector contributes approximately 18\% to the gross domestic product and employs nearly 42\% of the workforce \cite{fao2023}. Despite this immense economic and social significance, Indian farmers routinely face devastating price volatility alongside acute agronomic threats. Market commodity values can swing by 200\% to 400\% within a single season due to erratic monsoons and fragmented cold-chain infrastructure \cite{chand2012}. Furthermore, fungal, bacterial, and viral crop diseases compound these economic losses, severely handicapping viable crop yields prior to harvest.

A 2019 report by NITI Aayog estimated that market timing inefficiencies alone cost the Indian agricultural economy over \$8 billion annually in post-harvest losses \cite{nitiaayog2019}. Meanwhile, untimely diagnosis of foliar diseases decimates seasonal outputs. While existing government portals provide retrospective data, localized real-time predictive intelligence is strictly absent.

\subsection{Problem Statement}
This unified framework systematically resolves a triad of coupled agro-economic paradigms:
\begin{enumerate}
    \item \textbf{Crop Recommendation:} Given a specific localized soil profile (N, P, K, pH) and climatic baseline, classify the optimal agronomic crop.
    \item \textbf{Market Forecasting:} Predict the modal valuation of said commodities across a 14-day rolling horizon by extracting multidimensional weather shocks and arrival volume correlations.
    \item \textbf{Plant Pathology Diagnosis:} Ingest unconstrained RGB leaf imagery to instantly detect, identify, and categorize physiological plant stress states out of 38 distinctive biotic classes.
\end{enumerate}

\subsection{Literature Survey}
Presently, deep learning architectures—specifically Long Short-Term Memory (LSTM) networks \cite{hochreiter1997} and spatial-temporal Transformers \cite{vaswani2017}—are considered state-of-the-art for multivariate forecasting. Within this ecosystem, the Temporal Fusion Transformer (TFT) represents a breakthrough in heterogeneous time-series integration by individually gating static metadata and dynamic exogenous variables \cite{lim2021tft}. Furthermore, the evaluation of "zero-shot" models like Moirai against TFT establishes a novel comparative boundary \cite{woo2024moirai}.

In the domain of pathology, localized visual feature extraction via classical computer vision has been rapidly marginalized by Convolutional Neural Networks (CNNs). Top-tier efficient architectures like EfficientNet \cite{tan2019efficientnet} optimize model scaling linearly across depth, width, and resolution, presenting ideal deployment characteristics for resource-constrained edge systems typical in agricultural implementations.

\subsection{Objectives}
The core technical objectives include:
\begin{enumerate}
    \item Engineering highly associative tabular data spaces linking soil biometrics, NASA weather feeds, and Agmarknet metrics.
    \item Training and evaluating a Temporal Fusion Transformer (TFT) against gradient-boosted trees (LightGBM) for 14-day commodity regression.
    \item Executing transfer learning upon an EfficientNet-B0 backbone to yield an automated diagnostic classifier for plant diseases.
    \item Fusing these discrete inferential processes into an integrated, serverless-capable HTTP dashboard endpoint.
\end{enumerate}

\section{Methodology}

The project architecture comprises three decoupled machine learning backends bridged via a centralized web-hosted inference protocol.

\subsection{System 1: Market Price Forecasting}

\subsubsection{Data Aggregation \& Feature Engineering}
A custom scraper aggregated historical commodity pricing directly from Agmarknet endpoints. Using fuzzy heuristic matching, the data was enriched with location-specific daily indices obtained from the NASA POWER API (Temperature, Humidity, Irradiance, Rainfall).
Over 70 continuous and discrete features were engineered, preserving cyclic temporal boundaries via Fourier transformations, alongside rolling lag indicators specifying supply-chain volatility over 7-to-30 day intervals.

\subsubsection{Temporal Fusion Transformer (TFT) Setup}
While a gradient-boosted tree (LightGBM) generated a robust statistical baseline, complex interconnected shock parameters (sudden intense precipitation correlated with future low supply) necessitated the deployment of a Temporal Fusion Transformer via PyTorch Forecasting. 
The TFT explicitly modeled `Mandi` and `Commodity` as static categorical embeddings. Time-varying unknown reals such as Target Price, Volatility, and Rolling Averages were processed via customized Gated Residual Networks (GRNs). The network was optimized against a \texttt{QuantileLoss} metric across a distribution of 7 probability quantiles, yielding not just a point forecast, but an upper and lower economic confidence band representing volatility risk accurately.

\subsection{System 2: Crop Recommendation}
Operating via Kaggle's premier Soil Dataset, the 7 fundamental agronomic properties were expanded exponentially into 45 composite multi-variate features (e.g., $N \times P \times K$ multiplicative interactions, aridity deviations). A highly calibrated categorical LightGBM tree-ensemble was trained against these multi-dimensional boundaries, bypassing arbitrary non-linear scaling limitations intrinsic to standard logistic formulations.

\subsection{System 3: Plant Pathology Intelligence}

\subsubsection{Image Preprocessing \& Architecture}
To identify specific bacterial blights, rusts, and fungal aggregations, we instantiated a Convolutional Neural Network governed by the `EfficientNet-B0` morphology—selected for its maximum accuracy-to-parameter ratio. 
Leaf images strictly bounded to $(224 \times 224 \times 3)$ pixel densities were structurally normalized and augmented (via rotational and shearing tensor manipulations) representing variance in field-lighting and rotational capture.

\subsubsection{Transfer Learning}
Due to the vast structural complexity required for boundary-edge feature mapping, initialized base weights were adopted directly from \textit{ImageNet}. The spatial base layers were programmatically frozen (`trainable = False`), while a custom dense bottleneck (256 units, 30\% dropout) culminating in a 38-class softmax activation node was appended. The classification head was rapidly fine-tuned via the Adam optimizer ($lr=1e-3$) against the `categorical\_crossentropy` loss matrix utilizing 1500 images per batch across multi-fold validations.

\section{Quantitative Results}

\subsection{Forecasting Benchmark (TFT vs LightGBM vs Moirai)}
Extrapolating test sequences natively on unobserved 14-day horizons, the Temporal Fusion Transformer showcased robust superiority, specifically inside erratic anomaly patterns.

\begin{table}[htbp]
\caption{Global Test-Set Price Forecasting Performance}
\begin{center}
\begin{tabular}{lrr}
\toprule
\textbf{Model Architecture} & \textbf{SMAPE (\%)} & \textbf{Accuracy (\%)} \\
\midrule
Naive Baseline       & 48.40 & 51.60 \\
Moirai (Zero-Shot)   & 28.15 & 71.85 \\
LightGBM (Gradient Booting) & 9.69 & 90.31 \\
\textbf{Temporal Fusion Transformer (TFT)} & \textbf{7.82} & \textbf{92.18} \\
\bottomrule
\end{tabular}
\label{tab:price_metrics}
\end{center}
\end{table}

While LightGBM handled rigid cyclicity natively, TFT successfully minimized variance via dynamic temporal attention, pushing structural accuracy above 92.1\%. Foundation models (Moirai) failed to correlate effectively without the integration of contextual static variables representing the localized geo-spatial domains.

\subsection{Pathology Classification Performance}
Transfer learning optimizations executed on EfficientNet-B0 conclusively aligned probabilities. During out-of-sample inferencing, the customized CNN maintained an exceptional \textbf{98.4\% Macro F1-Score}. Inter-class categorical confusion was largely mitigated except in instances exhibiting compound biological infection (e.g., Early Blight overlapping concurrent Septoria leaf spot on Tomato subsets).
Similarly, the secondary LightGBM pipeline executed mapping of the 45-feature agronomy array with practically absolute precision (99.3\% validation convergence).

\section{Deployment \& System Unification}
In avoidance of heavyweight deployment silos commonly associated with Streamlit clusters or distributed Django modules, we authored an ultra-lightweight monolithic Python HTTP Server (`http.server`).
This architecture securely hosts the pre-trained `best\_model.keras`, the `lightgbm.pkl` matrix, and static historical trajectory binaries in local memory space at `localhost:8080`. API POST triggers securely route structured tabular metrics and sequential byte arrays toward the specific AI module—facilitating sub-100ms algorithmic turnover rendered beautifully over an integrated asynchronous CSS framework.

\section{Conclusion}
The \textbf{AgriSense AI} platform effectively solves a complex array of biological and economic liabilities plaguing modern localized agriculture. By transitioning fundamental tree-based models (LightGBM) to heavily monitored attention frameworks like the Temporal Fusion Transformer (TFT) for hyper-accurate volumetric price forecasting, and fusing it intrinsically with cutting-edge Deep Learning computer vision pipelines (EfficientNet-B0) dedicated to foliar pathology, this project establishes a definitive ecosystem for data-driven precision farming. 

\begin{thebibliography}{99}

\bibitem{fao2023}
FAO, \emph{The State of Food and Agriculture 2023}, Food and Agriculture Organization of the United Nations, 2023.

\bibitem{nitiaayog2019}
NITI Aayog, \emph{Demand and Supply Projections Towards 2033}, Government of India, 2019.

\bibitem{chand2012}
R. Chand, ``Development policies and agricultural markets,'' \emph{Economic and Political Weekly}, vol.~47, no.~52, pp.~53--63, 2012.

\bibitem{tan2019efficientnet}
M. Tan and Q. Le, ``EfficientNet: Rethinking model scaling for convolutional neural networks,'' in \emph{Proc. ICML}, 2019, pp.~6105--6114.

\bibitem{lim2021tft}
B. Lim \emph{et al.}, ``Temporal Fusion Transformers for interpretable multi-horizon time series forecasting,'' \emph{International Journal of Forecasting}, vol.~37, no.~4, pp.~1748--1764, 2021.

\bibitem{woo2024moirai}
G. Woo \emph{et al.}, ``Unified training of universal time series forecasting transformers,'' in \emph{Proc. ICML}, 2024.

\bibitem{hochreiter1997}
S. Hochreiter and J. Schmidhuber, ``Long short-term memory,'' \emph{Neural Computation}, vol.~9, no.~8, pp.~1735--1780, 1997.

\bibitem{vaswani2017}
A. Vaswani \emph{et al.}, ``Attention is all you need,'' in \emph{Proc. NeurIPS}, 2017.

\end{thebibliography}

\end{document}
